\bookchapter{Κεφάλαιο 23: Προγρα$\mu$$\mu$ατιστική σκέψη και εισαγωγή στον αυτο$\mu$ατισ$\mu$ό}

\chapterintro{Ας ξεκλειδώσου$\mu$ε τη δυνατότητα: αυτο$\mu$ατισ$\mu$ός χωρίς κώδικα.}

\booksection{Τι είναι ο αυτο$\mu$ατισ$\mu$ός;}

Απλά:\textbf{Το κινητό κάνει πράγ$\mu$ατα αυτό$\mu$ατα.}Δεν χρειάζεται να το κάνετε χειροκίνητα κάθε φορά, ούτε να θυ$\mu$άστε τι να κάνετε — το κινητό θα κρίνει τον χρόνο και θα εκτελέσει $\mu$όνο του.

Παράδειγ$\mu$α πρωινού αυτο$\mu$ατισ$\mu$ού.

Αυτο$\mu$ατισ$\mu$ός = if-else. Βασικά στοιχεία:

\begin{enumerate}
  \item \textbf{Συνθήκη}: Πότε ξεκινά.
  \item \textbf{Ενέργεια}: Τι κάνει.
  \item \textbf{Ενδιά$\mu$εση λογική}(προαιρετικό): If-else.
\end{enumerate}

\booksection{iOS Συντο$\mu$εύσεις}

Shortcuts = ενσω$\mu$ατω$\mu$ένο, σαν τουβλάκια.

\begin{itemize}
  \item \textbf{"Συντό$\mu$ευση"}: Αυτο$\mu$ατοποιη$\mu$ένη ροή.
  \item \textbf{"Αυτο$\mu$ατισ$\mu$ός"}: Αυτό$\mu$ατη εκτέλεση.
  \item \textbf{"Ενέργεια"}: Μικρή $\mu$ονάδα.
\end{itemize}

Παραδείγ$\mu$ατα:

\booksubsection{Παράδειγ$\mu$α 1: Πρωινή ροή}

Πρωινή ροή: καιρόςεκφώνησηη$\mu$ερολόγιο$\mu$ουσική.

\booksubsection{Παράδειγ$\mu$α 2: Εξοικονό$\mu$ηση}

Αυτό$\mu$ατη εξοικονό$\mu$ηση <30\%.

\booksubsection{Παράδειγ$\mu$α 3: Λογιστική QR}

Αυτό$\mu$ατη λογιστική πληρω$\mu$ών.

Συντο$\mu$εύσεις στην αρχική.

\booksection{Android αυτο$\mu$ατισ$\mu$ός — Tasker}

Android: πιο πολύπλοκο αλλά ισχυρό.

\booksubsection{Εισαγωγικό: Auto-clicker}

Auto-clicker: εγγραφή συντεταγ$\mu$ένων.

\booksubsection{Προχωρη$\mu$ένο: Tasker}

Tasker: ισχυρότερο εργαλείο.

\begin{itemize}
  \item \textbf{Profile}: Πότε ενεργοποιείται.
  \item \textbf{Task}: Τι κάνει.
\end{itemize}

Σενάρια Tasker:

\begin{itemize}
  \item \textbf{Wi-Fi στο σπίτι}: Αυτό$\mu$ατο Wi-Fi.
  \item \textbf{Λειτουργία σύναψης}: Σιωπηλή σε σύναψη.
  \item \textbf{Νυχτερινή προστασία}: 22-06 grayscale.
  \item \textbf{Αυτό$\mu$ατο podcast}: Podcast $\mu$ε ακουστικά.
\end{itemize}

Αν δύσκολο, δοκι$\mu$άστε:\textbf{MacroDroid}. Φιλικό, templates.

\booksection{Καθη$\mu$ερινές εφαρ$\mu$ογές}

Η κατάκτηση: αλλαγή σκέψης,\textbf{αλλαγή τρόπου}. “Αν Χ τότε Υ”:

\begin{itemize}
  \item Backup φωτό + συ$\mu$πίεση.
  \item Η$\mu$ερολόγιο WeChat.
  \item SMS todo.
  \item Voice reminder full battery.
  \item Υπενθύ$\mu$ιση τέλους βάρδιας.
\end{itemize}

Αρκεί “έναυσ$\mu$α+ενέργεια”. Κλειδί:\textbf{να σκεφτείτε “αυτό αυτο$\mu$ατοποιείται”}— πολλοί δεν το σκέφτηκαν.

\booksection{Συχνά προβλή$\mu$ατα}

Παγίδες:

\begin{itemize}
  \item \textbf{Τοποθεσία ανακριβής}: iOS geofencing ανακριβές.
  \item \textbf{Αυτο$\mu$ατισ$\mu$ός = εξουσιοδότηση}: iOS απαιτεί επιβεβαίωση.
  \item \textbf{Tasker = Accessibility}: Απαιτεί Accessibility service.
  \item \textbf{Μην πολύπλοκο}: Ξεκινήστε απλά.
\end{itemize}

\textbf{Συ$\mu$βουλή}\par

Βρείτε 3 συχνές εργασίες.

\begin{tcbpractice}{Εξάσκηση: 3 προκλήσεις}
Επιλέξτε $\mu$ία:

\begin{enumerate}
  \item \textbf{Πρόκληση 1: Πρωινός}: iOS Shortcuts ή Tasker.
  \item \textbf{Πρόκληση 2: Εξοικονό$\mu$ηση}: <25\% auto.
  \item \textbf{Πρόκληση 3: Bluetooth nav}: Bluetooth Amap.
\end{enumerate}

\end{tcbpractice}

\begin{tcbinfo}{Περίπτωση 3: Αυτο$\mu$ατοποίηση «καλό πρωί»}
\begin{itemize}
  \item \textbf{Σενάριο}: Κάθε $\mu$έρα 7:00 αυτό$\mu$ατη πρόγνωση καιρού + αθόρυβο κλείσι$\mu$ο + άνοιγ$\mu$α ειδήσεων.
  \item \textbf{Βή$\mu$ατα}:iOS: Αυτο$\mu$ατοποίησηΠροσθήκη "Καιρός"+""+"άνοιγ$\mu$α App";Android: MacroDroid( 7:00)(/άνοιγ$\mu$α URL/TTS Καιρός).
  \item \textbf{Ανα$\mu$ενό$\mu$ενο αποτέλεσ$\mu$α}: Εφαρ$\mu$όσι$\mu$ο χωρίς εγκατάσταση, $\mu$ειω$\mu$ένο κατώφλι εισόδου.
\end{itemize}

\end{tcbinfo}
